\subsection{Solving a System -- Elimination Method} 
The other common method for solving a linear system is the \emph{elimination} method. This method is harder to get our head around, but easier in most cases.
\begin{Example}
Use elimination to solve
\[
\begin{system}
3x-5y=21\\
4x+3y=-1
\end{system}
\]
\begin{lpic}{}{7}
\end{lpic}
The solution is \blankpair
\end{Example}
\begin{enumerate}
\item Put both equations in standard form. (We'll usually choose elimination when this is \makeblank, and we'll \makeblank[1in] the restrictions on $A$, $B$, and $C$.)
\item Multiply both sides of each equation by numbers chosen so that one variable will have opposite coefficients. (Swapping coefficients and switching one sign will always work.)
\item Add together the left-hand sides and the right-hand sides.
\item You now have an equation in \makeblank[1in] variable. Solve it.
\item You now have the \makeblank[1in] of one variable.
\item Plug into one of the equations to find the value of the \makeblank[1in] variable.
\item Check by plugging into the other equation.
\end{enumerate}